%% Independent caption file. Every caption is self-contained: readable without
%% the body text, without cross-references, and without float numbers.

\newcommand{\CapLocks}{%
The \NLocks{} obstacles standing between this project and the contact-gated
benchmark it intended to be evaluated on, as recorded by the project before any
computation was provisioned. The obstacles are sorted by what would remove them, and the
sort is the point: \NPermissionLocks{} of them are decisions by a person or an
institution, \NArtifactLocks{} is an evaluation file absent from the retained
workspace that
arrives only with one of those decisions, and \NSubstitutionLocks{} are refusals
the project imposed on itself so that a locally available corpus could not be
scored under the benchmark's name. No
obstacle in any of the three groups is removed by additional computation, by
additional labelling of the corpus already held, or by a larger model, which is
why the record was written before computation was provisioned.%
}

\newcommand{\CapPermissions}{%
The \NPermissionLocks{} permissions the project's record lists, each against the
\NPermissionLocks{} things a project needs before the benchmark can be run. A
cell reads \emph{delivers} only where holding the permission on that row
produces the item in that column, and \emph{still required} everywhere else. The
relation is exactly the diagonal: each approval yields one of the three and
leaves the remaining \NPermissionPairs{} cells untouched, so holding one is not
partial progress toward the others. The row labels are the approvals themselves,
which are decisions by a person or an institution; the column labels are the
things those decisions produce, which are a file, a personal credential and
access to the images. Nothing on this grid is moved by computation, by a larger
model, or by labelling more of the corpus already held.%
}

\newcommand{\CapRequired}{%
How many labelled reports each of the \NFindings{} findings would need
before a draw is expected to contain \MinPresent{} positives, computed from the
prevalence of that finding in the frozen test split and plotted on a
logarithmic axis. The upper dashed line marks the whole test split of \NTest{}
reports and the lower dashed line marks the \NCohort{} reports that were
actually labelled. Bars above the upper line, drawn in red, mark the
\NExceedSplit{} findings for which the entire test split is too small: for
those, labelling every remaining report in the split would still not reach the
gate, so the limitation belongs to the corpus rather than to the size of the
pilot. Bars below it mark findings for which a larger pilot within the same
split would suffice. The x-axis labels are compact finding codes; every bar is
one finding, not one report. The compact x-axis key is: \FindingKey%
}

% Figures 3--5 use short, stable codes so the x-axis remains readable at the
% printed width. The full schema names stay in the evidence-derived tables and
% this key makes each lifted figure self-contained.
\newcommand{\FindingKey}{%
\shortstack[l]{%
\texttt{F01}=Enlarged Cardiomediastinum; \texttt{F02}=Cardiomegaly;\\
\texttt{F03}=Lung Opacity; \texttt{F04}=Lung Lesion;\\
\texttt{F05}=Edema; \texttt{F06}=Consolidation;\\
\texttt{F07}=Pneumonia; \texttt{F08}=Atelectasis;\\
\texttt{F09}=Pneumothorax; \texttt{F10}=Pleural Effusion;\\
\texttt{F11}=Pleural Other; \texttt{F12}=Fracture;\\
\texttt{F13}=Support Devices; \texttt{F14}=No Finding.%
}%
}

\newcommand{\CapLabelStates}{%
Composition of the four states emitted by the released CheXbert labeller for
each of the \NFindings{} findings across the \NTest{} reports in the frozen
test split. Each horizontal bar is a complete partition of that split: blue is
an explicitly positive label, light blue an explicitly negative label, orange
an uncertain label, and grey is \emph{not mentioned (not negative)}. The printed
positive count and percentage are the only values annotated on the bars. These
are labeller-output counts, not clinician adjudication or clinical disease
prevalence; in particular, a grey segment must not be read as absence. The rows
are ordered by the positive count so the rare findings remain visible.%
}

\newcommand{\CapDraw}{%
Whether the pre-registered subsample of \NCohort{} reports resembles the
\NTest{}-report split it was drawn from, for each of the \NFindings{}
findings. Open circles give the number of positives the split's own
prevalence predicts for a draw of this size; filled circles give the number the
draw actually contains, and the connecting line is the discrepancy. The dashed
line marks the presence gate of \MinPresent{} gold positives below which the
analysis reports no rate. One finding, drawn in red, differs from the split
by more than sampling alone comfortably explains, with the exact two-sided
hypergeometric probability printed beside it; the comparison threshold after
Bonferroni correction across all \NFindings{} findings is \Bonferroni{}. The
selection rule was a salted hash fixed before the labels were read, so the
discrepancy is sampling variation in a pre-registered draw rather than a choice
made after seeing the data. Open and filled markers use the same finding order
as the compact x-axis labels; the gate is a count of gold positives, not a
percentage of the draw. The compact x-axis key is: \FindingKey%
}

\newcommand{\CapAssertions}{%
What the two controls assert against what the \NCohort{} labelled reports
actually contain, for each of the \NFindings{} findings, counted in reports.
Grey bars give the gold positives present in the draw. Red bars give the reports
the majority control was labelled positive for, and blue bars the same for the
indication-retrieval control. The majority control emits one constant report for
every row, so its bars can only be zero or the full \NCohort{}: it reaches the
full height for exactly one finding and is zero everywhere else, including
for the finding that denotes a study with nothing to report. That single red
bar is the whole behaviour of a baseline usually assumed to predict normality,
and it is an artifact of how the automatic labeller reads one fixed string
rather than a property of the corpus. All bars share the \NCohort{}-report
denominator; the grey series is the gold-positive count against which the two
control assertions are read. The compact x-axis key is: \FindingKey%
}

\newcommand{\CapScored}{%
The two findings that clear the presence gate of \MinPresent{} gold
positives in the \NCohort{}-report subsample, and how much of each the two
controls missed. Points give the share of gold positives the control failed to
assert and the vertical bars extend to the one-sided \(90\%\) Clopper--Pearson
upper limit on that share, which is the honest width at these counts rather than
a decoration. The majority control misses every positive in both cells, so its
point sits at one and its interval has nowhere to extend. The retrieval control
does better in both cells, and its upper limits still reach close enough to one
that near-total failure cannot be excluded from these data. Both cells hold
fewer than fifteen positives, and the intervals are wide for that reason and not
for any other. The horizontal rule at one is the logical maximum miss rate; a
point at one is complete miss, not an unobserved cell.%
}
